% ==============================================================================
% styles/environments.tex
% ==============================================================================

% استایل ۱: باکس کامل با هدر رنگی مجزا (برای تعاریف و قضایا)
\tcbset{
    solid_box/.style={
        enhanced,
        breakable,
        arc=5pt,          % شعاع نرم گوشه‌ها
        boxrule=1.5pt,    % ضخامت خط دور
        fonttitle=\bfseries,
        coltitle=white,   % رنگ متن هدر
        top=4mm, bottom=4mm, left=4mm, right=4mm,
        toptitle=3mm, bottomtitle=3mm,
    }
}

% استایل ۲: باکس مینیمال با خط کناری ضخیم (برای مثال‌ها و نکات)
\tcbset{
    side_box/.style={
        enhanced,
        breakable,
        arc=4pt,          % شعاع نرم گوشه‌ها
        boxrule=0pt,      % بدون خط دور پیش‌فرض
        rightrule=5pt,    % خط ضخیم سمت راست (مناسب برای متون فارسی راست‌چین)
        leftrule=0pt,
        fonttitle=\bfseries,
        top=4mm, bottom=4mm, left=4mm, right=4mm,
        attach title to upper, % چسباندن تایتل به متن
        after title={\par\smallskip}, % رفتن به خط بعد پس از عنوان
    }
}

% شمارنده مشترک برای همه محیط‌های ریاضی که وابسته به فصل است
\newcounter{mathbox}[chapter]
% تنظیم نمایش شماره به صورت X.Y (مثلاً 1.2)
\renewcommand{\themathbox}{\thechapter.\arabic{mathbox}}

% ------------------------------------------------------------------------------
% محیط‌های شماره‌دار (استایل solid_box)
% ------------------------------------------------------------------------------

% 1. قضیه
\NewTColorBox[use counter=mathbox]{theorem}{ o }{
    solid_box,
    colback=BoxBlue, colframe=BoxBlueFrame,
    title={\faLightbulb \quad قضیه \themathbox\IfValueT{#1}{ \textnormal{(#1)}}}
}

% 2. لم
\NewTColorBox[use counter=mathbox]{lemma}{ o }{
    solid_box,
    colback=BoxBlue, colframe=BoxBlueFrame,
    title={\faLink \quad لم \themathbox\IfValueT{#1}{ \textnormal{(#1)}}}
}

% 3. نتیجه
\NewTColorBox[use counter=mathbox]{corollary}{ o }{
    solid_box,
    colback=BoxBlue, colframe=BoxBlueFrame,
    title={\faProjectDiagram \quad نتیجه \themathbox\IfValueT{#1}{ \textnormal{(#1)}}}
}

% 4. تعریف
\NewTColorBox[use counter=mathbox]{definition}{ o }{
    solid_box,
    colback=BoxGreen, colframe=BoxGreenFrame,
    title={\faBook \quad تعریف \themathbox\IfValueT{#1}{ \textnormal{(#1)}}}
}

% ------------------------------------------------------------------------------
% محیط‌های بدون شماره (استایل side_box)
% ------------------------------------------------------------------------------

% 5. مثال
\NewTColorBox[use counter=mathbox]{example}{ o }{
    side_box,
    colback=BoxGray, colframe=BoxGrayFrame, coltitle=BoxGrayFrame!90!black,
    title={\faCogs \quad مثال \themathbox\IfValueT{#1}{ \textnormal{(#1)}}}
}

% 6. حل
\NewTColorBox{solution}{}{
    side_box,
    colback=BoxGray, colframe=BoxGrayFrame, coltitle=BoxGrayFrame!90!black,
    title={\faCheckCircle \quad حل}
}

% 7. نکته مهم
\NewTColorBox{note}{ o }{
    side_box,
    colback=BoxYellow, colframe=BoxYellowFrame, coltitle=BoxYellowFrame!90!black,
    title={\faStar \quad نکته مهم\IfValueT{#1}{ \textnormal{(#1)}}}
}

% 8. هشدار
\NewTColorBox{warning}{ o }{
    side_box,
    colback=BoxRed, colframe=BoxRedFrame, coltitle=BoxRedFrame!90!black,
    title={\faExclamationTriangle \quad هشدار\IfValueT{#1}{ \textnormal{(#1)}}}
}

% ------------------------------------------------------------------------------
% محیط اثبات (Proof)
% ------------------------------------------------------------------------------
% استفاده از پکیج amsthm ولی با استایل سفارشی
\renewcommand{\qedsymbol}{$\blacksquare$} % مربع توپر
\renewenvironment{proof}{{\bfseries \textcolor{BoxBlueFrame!90!black}{\faPenFancy \quad اثبات.}\quad}}{\hfill\qedsymbol\newline}
